\subsubsection*{推论 2.3.}
\underline{在上述假设和记号下, 假设 $x_0\in F$, 并且对每个不同于 $x_0$ 的 $x\in F$, 通过 $x$ 的 $f$ 的纤维在 $x$ 处有普通二次奇点. 则}
\[
\mu(X/S,x_0)=n(X/S,x_0).
\]
对 $x\ne x_0$, 命题 1.13 给出
\[
\mu(X/S,x)=n(X/S,x),
\]
然后只须应用 2.2.

\subsubsection*{定理 2.4.}
\underline{设 $S$ 为剩余域完美的等特征 Hensel 离散赋值环谱, $f:X\to S$ 为有限型平坦态射, $\bar x$ 为 $X$ 的几何点, 其在 $X$ 中的像为 $X_s$ 的闭点 $x$, $X/S$ 在该点处不光滑, 且该点在这样的点中孤立. 假设 $X$ 在 $x$ 处正则. 则 $X/S$ 在 $\bar x$ 处的消灭圈总数（1.8）等于 Milnor 数 $\mu(X/S,x)$.}

设 $S'$ 为由 $S$ 将剩余域 $k(s)$ 扩张到其可分闭包得到的 Hensel 离散赋值环谱, 置 $X'=X\times_S S'$. 2.4 对 $X/S$ 的断言等价于对 $X'/S'$ 的断言. 因为已经假设 $k(s)$ 完美, 所以可假设 $S$ 的剩余域代数闭.

对这样的 $S$, 设 $S'$ 为其完备化, 置 $X'=X\times_S S'$. $X/S$ 和 $X'/S'$ 在 $x$ 处的 Milnor 数相同, 且由 XIII, 2.4.2, 消灭圈总数也相同. 因此还可假设 $S$ 完备, 从而同构于 $k[[t]]$, 其中 $k$ 代数闭. 若 $X_s$ 在 $x$ 处的维数为 $n$, 则有:

\subsubsection*{命题 2.5.}
\underline{存在 $k$ 上纯维数 $n+1$ 的光滑射影簇 $Y$, 一个在 $Y$ 的有限点集 $F$ 之外光滑的态射 $f:Y\to\mathbb P^1_k$, 以及点 $y_0\in F$, 使得:}
\begin{enumerate}
\item[(a)] \underline{若 $y\in F$ 且 $y\ne y_0$, 则 $y$ 为通过 $y$ 的 $f$ 的纤维的普通二次奇点;}
\item[(b)] \underline{$f(y_0)=0$, 且若 $a:S\to\mathbb P^1_k$ 为将 $S$ 与 $\mathbb P^1_k$ 在 $0$ 处的完备化自然等同的态射, 则 $X$ 在 $x$ 处的 Hensel 化与 $a^*Y$ 在 $y_0$ 处的 Hensel 化之间存在 $S$-同构.}
\end{enumerate}

说明 2.5 蕴含 2.4. 由 2.3 与 2.5(a), $Y/\mathbb P^1_k$ 在 $y_0$ 处的 Milnor 数等于消灭圈总数.\footnote{原印本此处写作 $Y/T$, 但并未引入底 $T$; 命题 2.5 所定义的态射为 $Y\to\mathbb P^1_k$.}
如上, XIII, 2.4.2 对 $a^*Y/S$ 给出相同的结果, 因而也对 $(X/S,x)$ 给出相同结果; 由 2.5(b), 后者与 $(a^*Y/S,y_0)$ 在平展拓扑下局部同构.
